\section{Operating Systems, Programs, and Process Mechanics}

Before Python can be understood as a programming language, it is useful to understand what happens underneath any language when code is executed. A program does not run in isolation, directly on the physical machine, simply because a programmer wrote instructions in a source file. Between source code and execution there are several layers: executable files, operating-system loaders, process memory layouts, virtual address spaces, kernel scheduling, and hardware-supported resource isolation. These layers form the execution environment in which both C programs and Python programs eventually run.

This section uses the traditional C execution model as the reference point. C is useful here because it exposes the native path clearly: source code is transformed through preprocessing, compilation, assembly, and linking into a binary executable; the operating system loads that executable; the loader maps code and data into virtual memory; the process receives a stack, heap, registers, arguments, and resources; and the CPU begins executing machine instructions from the program's text segment. This path gives us a concrete model of what a native program is and how it becomes an active process.

The main distinction developed in this section is the distinction between a \emph{program} and a \emph{process}. A program is a stored blueprint: a file or collection of files containing instructions, data, and metadata. A process is the running instance created from that blueprint. The process has its own virtual address space, its own stack, heap, mapped code regions, open resources, and scheduling identity. This distinction is essential because Python code will later be shown not as a native executable in the ordinary C sense, but as input executed inside an already running interpreter process.

We will also examine the role of the operating system kernel. The kernel is the privileged layer that creates processes, maps memory, protects address spaces, handles traps and system calls, schedules CPU time, and mediates access to files, devices, and other resources. Even a C program compiled to machine code does not own the machine. It runs inside a boundary created and enforced by the operating system. Its pointers are process-virtual addresses, its files are kernel-managed resources, and its execution can be suspended and resumed by the scheduler.

The memory layout of a native process is another central theme. We will separate the text segment, data segment, BSS, heap, and stack. These regions explain where machine instructions live, where global and static variables are stored, where dynamic allocations are created, and how function calls produce stack frames. This is the model that C programmers usually approach directly through pointers, arrays, \texttt{malloc()}, \texttt{free()}, function calls, and local variables.

Finally, we will connect process mechanics to multitasking. Modern operating systems do not run only one program until completion. They time-slice CPU resources, suspend and resume processes or threads, switch execution contexts, and distribute work across multiple cores. This explains the difference between processes and threads, between isolation and shared memory, and between operating-system scheduling and runtime-level scheduling. That last transition prepares the ground for later Python concepts such as generators, coroutines, event loops, and asynchronous execution.

The purpose of this section is not to turn the chapter into an operating-systems course. The purpose is to establish the execution model that Python will modify, abstract, and partially hide. Once the native C model is clear, the Python model becomes easier to describe precisely: the operating system runs the CPython executable as a normal native process, while the user's Python source file is parsed, compiled to bytecode, and executed inside that process by the Python runtime.
